\SourceLeaf{12}
Punkte ist also gerade oder ungerade, je nachdem, ob $M$ außerhalb oder innerhalb des festen Teils des Körpers liegt. Offenbar ist der Fall, in dem $M$ auf der Oberfläche des Körpers selbst liegt, dem ersten oder dem zweiten Fall zuzurechnen, je nachdem, ob der Radius $M\Pi$ sich anfangs vom festen Teil des Körpers entfernt oder in ihn eintritt. Man denke sich ferner Geraden von $M$ zum Rand des kleinen Flächenstücks $\mathrm{d}\Sigma$ gezogen. Sie bilden eine Kegelfläche (im weiteren Sinne) und begrenzen auf der Oberfläche unseres Körpers an den Punkten $P',P'',P'''$ usw. der Reihe nach die kleinen Flächenstücke $\mathrm{d}s',\mathrm{d}s'',\mathrm{d}s'''$ usw. Schließlich beschreibe man durch die Punkte $P',P'',P'''$ usw. kleine Stücke von Kugeloberflächen um den Mittelpunkt $M$ mit den Radien
\[
MP'=r',\quad MP''=r'',\quad MP'''=r'''\text{ etc.}
\]
{\rightskip=0pt plus 3em Die kleinen Flächenstücke, die der Kegel daraus ausschneidet, seien $\mathrm{d}\sigma',\mathrm{d}\sigma'',\mathrm{d}\sigma'''$ usw. Wir betrachten alle diese kleinen Flächenstücke $\mathrm{d}\Sigma,\mathrm{d}s',\mathrm{d}\sigma'$ usw. als positiv. Damit haben wir
\[
\mathrm{d}\Sigma=\frac{\mathrm{d}\sigma'}{r'r'}=\frac{\mathrm{d}\sigma''}{r''r''}=\frac{\mathrm{d}\sigma'''}{r'''r'''}\text{ etc.}
\]
Das kleine Flächenstück $\mathrm{d}\sigma'$ lässt sich als Projektion des kleinen Flächenstücks $\mathrm{d}s'$ auf eine Ebene betrachten, auf der die Gerade $P'M$ senkrecht steht. Daher gilt $\mathrm{d}\sigma'=\pm\mathrm{d}s'.\cos MQ'$, wobei das obere oder das untere Vorzeichen zu nehmen ist, je nachdem, ob $MQ'$ spitz oder stumpf ist. Der erste Fall tritt ein, wenn die von $P'$ nach $M$ gezogene Gerade sich vom Körper entfernt, also wenn $M$ außerhalb des Körpers liegt. Der zweite tritt ein, wenn die Gerade $P'M$ bei $P'$ in den Körper eintritt, also wenn $M$ innerhalb des Körpers liegt. Entsprechend gilt $\mathrm{d}\sigma''=\mp\mathrm{d}s''\cos MQ''$, $\mathrm{d}\sigma'''=\pm\mathrm{d}s'''.\cos MQ'''$ usw. Daraus folgt
\par}
I.\quad Liegt $M$ außerhalb des Körpers, so gilt
\[
\begin{aligned}
\mathrm{d}s'.\cos MQ'&=+r'r'\mathrm{d}\Sigma\\
\mathrm{d}s''.\cos MQ''&=-r''r''\mathrm{d}\Sigma\\
\mathrm{d}s'''.\cos MQ'''&=+r'''r'''\mathrm{d}\Sigma\\
&\text{etc.}
\end{aligned}
\]
\par
II.\quad Liegt $M$ dagegen innerhalb des Körpers, so gilt
\[
\begin{aligned}
\mathrm{d}s'.\cos MQ'&=-r'r'\mathrm{d}\Sigma\\
\mathrm{d}s''.\cos MQ''&=+r''r''\mathrm{d}\Sigma\\
\mathrm{d}s'''.\cos MQ'''&=-r'''r'''\mathrm{d}\Sigma\\
&\text{etc.}
\end{aligned}
\]
\par
Im Fall I gilt also (weil die Zahl der Gleichungen gerade ist)
